\begin{abstract}
  Realistic image synthesis is a central problem in Computer Graphics, aiming
  to algorithmically reproduce light transport in a physically plausible
  manner. \emph{Path Tracing} is one of the approaches that makes the fewest
  simplifications to this transport, allowing effects such as reflections and
  indirect illumination to arise naturally from the algorithm. However, since
  it estimates the rendering equation \citep{kajiya86} via a Monte Carlo
  integrator, the method is computationally expensive, and its quality depends
  directly on the number of samples used. In this work, we implement a
  \emph{Path Tracer} and explore Multiple Importance Sampling \citep{veach95},
  a technique that combines different sampling strategies to reduce estimator
  variance. The implementation combines BRDF sampling and direct light source
  sampling, and the derivation of the MIS weights is presented explicitly for
  the full path. The effectiveness of the technique is evaluated by comparing
  images rendered with an equal number of samples, with and without Multiple
  Importance Sampling. The results demonstrate that the technique significantly
  reduces estimator variance: variance maps rendered with an equal number of
  samples show an expressive reduction, with the version without Multiple
  Importance Sampling exhibiting visibly high variance in contrast to the
  version that employs it.

\end{abstract}
